\subsection[Tate bar-cobar admissibility and the infinity-categorical comparison]{Tate bar-cobar admissibility and the $\infty$-categorical
comparison of the cochain bar-cobar identity}\label{sec:tate-quillen}

Inside an admissible Tate envelope satisfying
Statement~\ref{stmt:tate-model-envelope}, the bar-cobar adjunction
\(\Omega\dashv\mathrm{Bar}\) is required to satisfy the admissible
unit and counit weak-equivalence hypotheses
(Theorem~\ref{thm:tate-bar-cobar-quillen}).  The CE/PV map
\(c_f\mapsto\theta_f,\ u_f\mapsto J(f)\) then gives an admissible
filtered equivalence between the closed CE cochain algebra and the
Schouten polyvector algebra of the Hamiltonian formal disk.  This is a
Koszul comparison for the trace sector.  It is not the construction of
the full Hochschild cochain complex, nor of the derived chiral centre;
those require the Hochschild-centre lift and its centrality homotopies.
After restriction to a holomorphic curve, the \(d=1\) matching to the
bulk side is the chiral-algebra theorem of
\cite{beilinson-drinfeld-chiral}.  The arity-two holomorphic collision
is the two-dimensional diagonal singular product of
Theorem~\ref{thm:formal-two-dimensional-holomorphic-ope}; the Tate
bar-cobar comparison controls only the Koszul algebraic coordinate of
that stalk:
\(\Omega C_*^{\mathrm{CE}}(\mathfrak g)\xrightarrow{\sim}\widehat
U(\mathfrak g)\) on the Lie side, dualized through exact continuous
duality to
\[
  C^\bullet_{\mathrm{CE},\mathrm{cont}}(\mathfrak g)
    \simeq A_{\mathrm{op}}^!,
\]
with the linear symbol \(B_f\mapsto O_f,\ \Theta_{\rho_{a,b}}\mapsto
\theta^{a,b}\) of
Definition~\ref{defn:universal-open-side-twisting-cochain}.

For zero-differential finite-dimensional pieces $\{V_d\}_{d\geq1}$,
the inclusion
\[
  V_\oplus=\bigoplus_{d\geq1}V_d
    \hookrightarrow
  V_\Pi=\prod_{d\geq1}V_d,
\]
filtered by degree tails, is an isomorphism on every finite quotient
and on associated graded, yet $H^0(V_\Pi/V_\oplus)=V_\Pi/V_\oplus\neq0$
(Lemma~\ref{lem:tate-envelope-tail-boundary}).  The tail-quotient class
is invisible to admissible filtered weak equivalence and constitutes
one summand of the Koszul obstruction datum (the five-summand
obstruction complex collecting the window, topology, conilpotence,
Maurer--Cartan, and cobar obstruction components for the twist \(\tau\)):
$\mathcal O_{\mathrm{Kos}}=\mathcal O_{\mathrm{win}}\oplus
\mathcal O_{\mathrm{top}}\oplus\mathcal O_{\mathrm{conilp}}\oplus
\mathcal O_{\mathrm{MC}}(\tau)\oplus\mathcal O_{\mathrm{cobar}}(\tau)$,
of Definition~\ref{defn:universal-koszul-obstruction-complex}.  The
recognition criterion
\[
  \omega_{\mathrm{win}}=0,\quad
  \mathcal O_{\mathrm{conilp}}=0,\quad
  \mathcal O_{\mathrm{MC}}(\tau)=0,\quad
  H^\bullet\bigl(\mathcal O_{\mathrm{top}}
    \oplus\mathcal O_{\mathrm{cobar}}(\tau)\bigr)=0
\]
is necessary and sufficient for filtered-cobar admissibility
(Theorem~\ref{thm:universal-filtered-cobar-recognition}).  It is not,
by itself, an ordinary cohomology equivalence of the chiral Hochschild
target, since Lemma~\ref{lem:tate-envelope-tail-boundary} exhibits
tail cohomology invisible to admissible weak equivalence.  An ordinary
Hochschild-cohomology statement requires a tail-vanishing or
Roos-acyclicity hypothesis.
Theorem~\ref{thm:tate-bar-cobar-quillen} descends to an equivalence of
associated \(\infty\)-categories
(Corollary~\ref{cor:tate-bar-cobar-infinity}).  The Weiss-cosheaf
condition on the unreduced factorization algebra over $\R^2\times\C^2$
remains separate (Proposition~\ref{prop:weiss-cosheaf-residual}); it
requires the brane-defect heat-kernel propagator extension of
Problem~\ref{prob:formal-factorization-center}~item~(2).

\subsubsection{The Tate chain category}

Let $\Cat{TateVec}$ be the closed symmetric monoidal category of
complete filtered Tate $\C$-vector spaces of \alpharef. The
filtration on every object is taken \emph{complete and Hausdorff}:
$V=\varprojlim_w V_{\leq w}$ and $\bigcap_w F^w V=0$. Concrete
morphisms are filtration-preserving continuous chain maps. The
working ambient category is the bounded-below filtered subcategory
$\Cat{TateChain}_{\geq 0}$ of non-negatively filtered chain
complexes in $\Cat{TateVec}$, large enough to contain every CE chain
coalgebra and completed enveloping algebra used here.  The exact
finite-window Mittag-Leffler category is constructed first.  The
model-categorical statements add transfer, smallness, and monoid
axioms to this exact category; their scope is the exact category
constructed here.

\begin{defn}[Strict finite-window Mittag-Leffler Tate category]
\label{defn:strict-ml-tate-category}
  Let \(\Cat{MLTateChain}_{\geq0}^{\mathrm{str}}\) be the category whose
  objects are inverse systems
  \[
    V=(V_{\leq M},\pi^V_{M,N})_{M\geq N\geq1}
  \]
  of bounded-below finite-dimensional chain complexes over \(\C\), with
  degreewise surjective transition maps.  The represented complete
  object is
  \[
    |V|=\varprojlim_MV_{\leq M},\qquad
    F^M|V|=\ker(|V|\to V_{\leq M-1}).
  \]
  Morphisms are compatible families of chain maps
  \(f_{\leq M}\colon V_{\leq M}\to W_{\leq M}\).  The completed tensor
  product is defined on paired windows by
  \[
    V\widehat\otimes_{\mathrm{ML}}W
      =
    \left(
      V_{\leq M}\otimes W_{\leq M}
    \right)_{M\geq1},
  \]
  with the induced surjective transition maps, and
  \[
    |V\widehat\otimes_{\mathrm{ML}}W|
      =
    \varprojlim_M(V_{\leq M}\otimes W_{\leq M}).
  \]
  The continuous dual attached to this strict tower is
  \[
    V^\vee_{\mathrm{ML}}=\varinjlim_MV_{\leq M}^\vee .
  \]
  A morphism is an admissible filtered weak equivalence when each
  \(f_{\leq M}\) and the induced associated-graded map are
  quasi-isomorphisms.

  The only universal property asserted for this category is the ordinary
  inverse-limit one: a compatible family \(x_M\in V_{\leq M}\) determines
  a unique element \(x\in |V|\), and every map out of a test strict tower
  is exactly a compatible family of finite-window maps.  Initial and
  terminal properties among all completions of a Tate vector space lie
  outside this assertion.
\end{defn}

\begin{prop}[Exactness and kernels in the strict ML category]
\label{prop:strict-ml-category-exactness}
  In \(\Cat{MLTateChain}_{\geq0}^{\mathrm{str}}\):
  \begin{enumerate}[label=(\roman*)]
    \item the filtration on \(|V|\) is complete and Hausdorff;
    \item inverse limits of strict surjective towers are exact, so
      \(H^\bullet(|V|)=\varprojlim_MH^\bullet(V_{\leq M})\) and
      \(\varprojlim^1_MH^{\bullet-1}(V_{\leq M})=0\);
    \item completed tensor products and continuous duals are computed by
      the displayed finite-window formulas;
    \item if \(H=(H_{\leq M})\) and \(D=(D_{\leq M})\) are paired strict
      ML systems and finite tensors
      \(C_M\in H_{\leq M}\otimes D_{\leq M}\) are compatible under
      transition, then
      \[
        C=\varprojlim_MC_M
          \in |H|\widehat\otimes_{\mathrm{ML}}|D|
      \]
      exists and is the unique strict ML kernel with those finite-window
      projections.
  \end{enumerate}
\end{prop}

\begin{proof}
  Completeness is the definition of \(|V|\) as an inverse limit; the
  intersection of the kernels \(F^M|V|\) is zero because an element whose
  image vanishes in every finite quotient is the zero compatible family.
  Surjectivity of the transition maps gives the Mittag-Leffler condition.
  The Milnor exact sequence then has zero \(\varprojlim^1\)-term, proving
  exactness and the cohomology formula.  The tensor and dual formulas are
  definitions inside the strict category and are finite-dimensional on each
  window.  Finally, a compatible tensor family \((C_M)_M\) is precisely an
  element of the inverse limit
  \(\varprojlim_M(H_{\leq M}\otimes D_{\leq M})\), which is the completed
  tensor product by definition.  The uniqueness is the inverse-limit
  universal property of Definition~\ref{defn:strict-ml-tate-category}.
\end{proof}

\begin{stmt}[Admissible Tate model envelope]\label{stmt:tate-model-envelope}
  An admissible Tate model envelope is a bounded-below, locally
  presentable symmetric monoidal model category
  \(\Cat{TateChain}_{\geq0}^{\mathrm{adm}}\) equipped with a fully
  faithful exact symmetric-monoidal embedding of the strict category
  \[
    \Cat{MLTateChain}_{\geq0}^{\mathrm{str}}
      \hookrightarrow
    \Cat{TateChain}_{\geq0}^{\mathrm{adm}},
  \]
  and containing the CE chains, completed enveloping algebras, completed
  symmetric algebras, and finite-window inverse systems used below.
  Weak equivalences are admissible filtered
  quasi-isomorphisms: quasi-isomorphisms on every finite quotient and on
  associated graded.  The tensor product restricts to
  \(\widehat\otimes_{\mathrm{ML}}\) on strict Mittag-Leffler systems, and
  filtered inverse limits with surjective transition maps are exact by
  Proposition~\ref{prop:strict-ml-category-exactness}.  The additional
  hypotheses are transfer and smallness for the bar-cobar adjunction, the
  monoid axiom for completed tensor products, and compatibility of
  \(\widehat\otimes\) with the finite-window model structures.  All
  Quillen statements below are relative to such an envelope.  The strict
  Mittag-Leffler category gives exact inverse limits; it does not by
  itself give local presentability, Hovey recognition, the monoid axiom,
  or operadic transfer.
\end{stmt}

\begin{lem}[Mittag-Leffler dictionary for closed Tate windows]
\label{lem:tate-mittag-leffler-dictionary}
  Let
  \[
    \{V_{\leq M}\}_{M\geq1}
    \quad\text{and}\quad
    \{W_{\leq M}\}_{M\geq1}
  \]
  be inverse systems of finite-dimensional complexes with surjective
  transition maps, and put
  \[
    V=\varprojlim_MV_{\leq M},\qquad
    W=\varprojlim_MW_{\leq M}.
  \]
  In the strict ML category, and hence in any admissible envelope
  containing it, the following assertions hold:
  \begin{enumerate}[label=(\roman*)]
    \item \(\varprojlim\) is exact on these systems, so
      \(H^\bullet(V)=\varprojlim_MH^\bullet(V_{\leq M})\);
    \item continuous dualization changes the inverse system into the
      direct system of finite duals:
      \(V^\vee_{\mathrm{cont}}=\varinjlim_MV_{\leq M}^\vee\);
    \item completed tensor products are computed windowwise:
      \(V\widehat\otimes W
        =\varprojlim_M(V_{\leq M}\otimes W_{\leq M})\) whenever the
      transition maps are the paired finite-window truncations;
    \item a map \(f\colon V\to W\) is an admissible filtered
      quasi-isomorphism iff each finite-window map
      \(f_{\leq M}\colon V_{\leq M}\to W_{\leq M}\) is a
      quasi-isomorphism and the induced associated-graded maps are
      quasi-isomorphisms.
  \end{enumerate}
\end{lem}

\begin{proof}
  Surjectivity of the transition maps gives the Mittag-Leffler
  condition, hence \(\varprojlim^1=0\).  The cohomology assertion is
  the standard Milnor exact sequence with the \(\varprojlim^1\)-term
  killed.  Continuous linear functionals on \(V\) factor through a
  finite quotient \(V_{\leq M}\), giving the displayed direct limit of
  duals.  The completed tensor product in
  Definition~\ref{defn:strict-ml-tate-category} is defined by the same
  finite-window inverse system, so item (iii) is tautological inside the
  strict category.  Item (iv) is the definition of weak equivalence in
  Definition~\ref{defn:strict-ml-tate-category} and
  Statement~\ref{stmt:tate-model-envelope}.
\end{proof}

\begin{defn}[Closed-window obstruction terms]
\label{defn:closed-window-obstruction-terms}
  Let \(\mathfrak h=\varprojlim_M\mathfrak h_{\leq M}\) be a complete
  filtered Lie algebra with finite-dimensional windows and surjective
  transition maps \(\pi_{M,N}\).  Write
  \(D_{\leq M}\) for a finite dual module paired with
  \(\mathfrak h_{\leq M}\), with surjective transition maps
  \(\pi^\vee_{M,N}\), and put
  \(D_{\mathrm{cont}}=\varprojlim_MD_{\leq M}\).  For \(M\geq N\)
  define:
  \[
    \mathfrak o^{\mathrm{br}}_{M,N}(x,y)
      =\pi_{M,N}[x,y]_M-[\pi_{M,N}x,\pi_{M,N}y]_N,
  \]
  \[
    \mathfrak o^{\mathrm{coad}}_{M,N}(x,\lambda)
      =\pi^\vee_{M,N}(x\cdot_M\lambda)
       -(\pi_{M,N}x)\cdot_N(\pi^\vee_{M,N}\lambda),
  \]
  and
  \[
    \mathfrak o^{\mathrm{Cas}}_{M,N}
      =(\pi_{M,N}\otimes\pi^\vee_{M,N})C_M-C_N .
  \]
  The closed window system is \emph{cochain-CE/PV-admissible} if
  \(\mathfrak o^{\mathrm{br}}\) and \(\mathfrak o^{\mathrm{coad}}\)
  vanish for all \(M\geq N\).  It is
  \emph{kernel-admissible} only after the additional Casimir term
  \(\mathfrak o^{\mathrm{Cas}}\) vanishes in the chosen completed tensor
  category.  The coordinate CE/PV cochain identity contains the bracket
  and coadjoint transition equations; kernel-admissibility adds the
  Casimir equation.
\end{defn}

\begin{lem}[Closed CE/PV finite-window compatibility and variance boundary]
\label{lem:closed-ce-pv-window-limit}
  Assume the closed window system of
  Definition~\ref{defn:closed-window-obstruction-terms} is
  cochain-CE/PV-admissible and each finite window has the finite cotangent
  CE/PV identity
  \[
    \Phi_M^{\mathrm{CE/PV}}\colon
    C^\bullet_{\mathrm{CE}}\bigl(\mathfrak h_{\leq M}\ltimes
      D_{\leq M}[1]\bigr)
    \xrightarrow{\sim}
    \mathrm{PV}\bigl(\mathbf S(\mathfrak h_{\leq M})\bigr).
  \]
  Then the \(\Phi_M^{\mathrm{CE/PV}}\) are compatible as finite coordinate tests under the
  contravariant CE pullbacks and the corresponding PV coordinate maps.
  These maps form compatible finite coordinate tests.  A completed
  admissible filtered quasi-isomorphism
  \[
    \Phi^{\mathrm{CE/PV}}_{\mathrm{cont}}\colon
    C^\bullet_{\mathrm{CE},\mathrm{cont}}
      \bigl(\mathfrak h\ltimes D_{\mathrm{cont}}[1]\bigr)
    \xrightarrow{\sim}
    \mathrm{PV}\bigl(\widehat{\mathbf S}_{\mathrm{Tate}}(\mathfrak h)\bigr)
  \]
  follows only after one gives either a chain-side inverse limit with
  exact continuous dualization or a topology whose cochain transition maps
  are dg maps.  If the kernel-admissible Casimir condition is also
  given, the same finite windows carry compatible \(P_0\)-pairings and
  kernels in that chosen tensor category.
\end{lem}

\begin{proof}
  The finite coordinate map \(\Phi_M^{\mathrm{CE/PV}}\) is determined on generators by
  \(c\mapsto\theta\) and \(u\mapsto O\).  Vanishing of
  \(\mathfrak o^{\mathrm{br}}_{M,N}\) says the CE differential on the
  Hamiltonian ghosts commutes with transition maps.  Vanishing of
  \(\mathfrak o^{\mathrm{coad}}_{M,N}\) says the cotangent-coordinate
  part of the CE differential, and equivalently the Lichnerowicz
  differential on \(O\)-coordinates, commutes with transition maps.
  Therefore the finite squares with
  \(\Phi_M^{\mathrm{CE/PV}}\) and \(\Phi_N^{\mathrm{CE/PV}}\) commute in
  the variance appropriate to cochains.  A quotient
  \(\mathfrak h_{\leq M}\twoheadrightarrow\mathfrak h_{\leq N}\) induces
  CE cochain maps by pullback along the quotient.
  Thus inverse-limit language on cochains requires the additional
  continuous-dualization or topology hypothesis stated above.  Vanishing
  of \(\mathfrak o^{\mathrm{Cas}}_{M,N}\) is the separate statement that
  the finite \(P_0\)-brackets and BV coefficient pairings are compatible
  under transition in the chosen kernel-admissible tensor category.
\end{proof}

\subsubsection{Twisting cochains and the filtered-cobar admissibility criterion}

\begin{defn}[Open-side twisting cochain datum]
\label{defn:universal-open-side-twisting-cochain}
  Fix an admissible filtered-cobar datum
  \[
    \mathfrak h=\varprojlim_M\mathfrak h_{\leq M},\qquad
    D=\varinjlim_M\mathfrak h_{\leq M}^{\vee},\qquad
    \mathfrak g=\mathfrak h\ltimes D[1],
  \]
  a complete augmented filtered open \(A_\infty\)-algebra
  \(A_{\mathrm{op}}\), and the completed CE/PV target
  \[
    B_{\mathrm{cl}}
      =
      C^\bullet_{\mathrm{CE},\mathrm{cont}}(\mathfrak g)
      \cong
      \mathrm{PV}\bigl(\widehat{\mathbf S}_{\mathrm{coord}}(\mathfrak h)\bigr).
  \]
  Put \(C_{\mathrm{op}}=\mathrm{Bar}(A_{\mathrm{op}})\), with reduced bar
  differential \(b_{\mathrm{op}}\), coproduct \(\Delta\), and iterated
  reduced coproducts \(\Delta^{(r)}\).  The bar coalgebra is required to
  be locally conilpotent in the completed sense: for every finite codomain
  window of \(B_{\mathrm{cl}}\) and every finite coaugmentation quotient of
  \(C_{\mathrm{op}}\), the iterated reduced coproducts \(\Delta^{(r)}\)
  vanish for all sufficiently large \(r\).  Equivalently, the
  conilpotence obstruction in
  Definition~\ref{defn:universal-koszul-obstruction-complex} is zero.
  This condition makes the Maurer--Cartan curvature below a continuous
  map and gives the twisting-cochain datum.

  An open-side twisting cochain is a continuous degree-one map
  \[
    \tau\colon C_{\mathrm{op}}\longrightarrow B_{\mathrm{cl}}
  \]
  vanishing on the coaugmentation and satisfying the Maurer--Cartan
  equation
  \[
    d_{B_{\mathrm{cl}}}\tau+\tau b_{\mathrm{op}}
      +\sum_{r\geq2}
        m_r^{B_{\mathrm{cl}}}\bigl(\tau^{\otimes r}\bigr)
        \Delta^{(r)}
      =0 .
  \]
  For the dg \(P_0\)-algebra \(B_{\mathrm{cl}}\), the higher
  multiplications vanish and this reads
  \[
    d_{B_{\mathrm{cl}}}\tau+\tau b_{\mathrm{op}}
      +\mu_{B_{\mathrm{cl}}}(\tau\otimes\tau)\Delta=0 .
  \]
  Equivalently,
  \[
    d_{\mathrm{conv}}\tau+\frac12[\tau,\tau]_{\mathrm{conv}}=0,\qquad
    d_{\mathrm{conv}}f=d_{B_{\mathrm{cl}}}f
      -(-1)^{|f|}f\,b_{\mathrm{op}},
  \]
  in the continuous convolution dg Lie algebra
  \(\Hom_{\mathrm{cont}}(C_{\mathrm{op}},B_{\mathrm{cl}})\).

  The twisting cochain induces the completed cobar morphism
  \[
    \kappa_\tau\colon
    \Omega C_{\mathrm{op}}
    \longrightarrow B_{\mathrm{cl}}.
  \]
  Its associated-graded symbol is the map
  \[
    \operatorname{gr}\kappa_\tau\colon
    \operatorname{gr}\Omega\mathrm{Bar}(A_{\mathrm{op}})
    \longrightarrow
    \operatorname{gr}B_{\mathrm{cl}}.
  \]
  In the Hamiltonian trace presentation the required linear symbol is
  \[
    \operatorname{gr}\kappa_\tau(B_f)=O_f,\qquad
    \operatorname{gr}\kappa_\tau(\Theta_{\rho_{a,b}})=\theta^{a,b},
  \]
  after the scalar trace is removed and the cotangent generator is read
  in the principal-part local-dual basis.
\end{defn}

\begin{defn}[Open-side Koszul obstruction complex]
\label{defn:universal-koszul-obstruction-complex}
  In the notation of
  Definition~\ref{defn:universal-open-side-twisting-cochain}, put
  \(D_{\leq M}=\mathfrak h_{\leq M}^\vee\),
  \(\mathrm{PV}_{\leq M}=\mathrm{PV}(\mathbf S(\mathfrak h_{\leq M}))\),
  and \(B_{\mathrm{cl},\leq M}\) for the corresponding finite CE/PV
  window.  The obstruction datum is
  \[
    \mathcal O_{\mathrm{Kos}}(\mathfrak h,D,A_{\mathrm{op}},\tau)
      =
      \mathcal O_{\mathrm{win}}
      \oplus
      \mathcal O_{\mathrm{top}}
      \oplus
      \mathcal O_{\mathrm{conilp}}
      \oplus
      \mathcal O_{\mathrm{MC}}(\tau)
      \oplus
      \mathcal O_{\mathrm{cobar}}(\tau).
  \]
  The window and conilpotence summands are degree-zero:
  \[
    \mathcal O_{\mathrm{win}}
      =
      \prod_{M\geq N}
      \left(
        \Hom(\Lambda^2\mathfrak h_{\leq M},\mathfrak h_{\leq N})
        \oplus
        \Hom(\mathfrak h_{\leq M}\otimes D_{\leq M},D_{\leq N})
      \right)
      \oplus
      Q_{\mathrm{Cas}},
  \]
  where
  \[
    Q_{\mathrm{Cas}}
      =
      \frac{\varprojlim_M(\mathfrak h_{\leq M}\otimes D_{\leq M})}
           {\operatorname{im}(\mathfrak h\widehat\otimes D)} .
  \]
  The distinguished window class is
  \[
    \omega_{\mathrm{win}}
      =
      \bigl(
        \mathfrak o^{\mathrm{br}}_{M,N},
        \mathfrak o^{\mathrm{coad}}_{M,N},
        [(C_{\leq M})_M]
      \bigr),
  \]
  where the first two terms are those of
  Definition~\ref{defn:closed-window-obstruction-terms} and the last
  term is the compatible finite-window Casimir family modulo strict
  completed tensors.  The topology summand is
  \[
    \mathcal O_{\mathrm{top}}^n
      =
      \varprojlim\nolimits^1_M
        H^{n-1}\!\left(C^\bullet_{\mathrm{CE}}(\mathfrak g_{\leq M})\right)
      \oplus
      \varprojlim\nolimits^1_M
        H^{n-1}\!\left(\mathrm{PV}_{\leq M}\right).
  \]
  The conilpotence summand is
  \[
    \mathcal O_{\mathrm{conilp}}
      =
      \left(
        C_*^{\mathrm{CE}}(\mathfrak g)/
        \bigcup_rF^{(r)}_{\mathrm{coaug}}
          C_*^{\mathrm{CE}}(\mathfrak g)
      \right)
      \oplus
      \bigcap_{r\geq1}\mathfrak g^{(r)}.
  \]

  The Maurer--Cartan summand is the explicit curvature
  \[
    \mathcal O_{\mathrm{MC}}(\tau)
      =
      \operatorname{MC}(\tau)
      =
      d_{\mathrm{conv}}\tau+\frac12[\tau,\tau]_{\mathrm{conv}}
      \in
      \Hom^2_{\mathrm{cont}}(C_{\mathrm{op}},B_{\mathrm{cl}}).
  \]
  The cobar obstruction complex is
  \[
    \mathcal O_{\mathrm{cobar}}(\tau)
      =
      \operatorname{Cone}(\operatorname{gr}\kappa_\tau)
      \oplus
      \bigoplus_n
      \varprojlim\nolimits^1_M
        H^{n-1}\!\left((\Omega C_{\mathrm{op}})_{\leq M}\right)[-n]
      \oplus
      \bigoplus_n
      \varprojlim\nolimits^1_M
        H^{n-1}\!\left(B_{\mathrm{cl},\leq M}\right)[-n].
  \]
\end{defn}

\begin{thm}[Filtered-cobar admissibility criterion]
\label{thm:universal-filtered-cobar-recognition}
  Fix a datum as in
  Definitions~\ref{defn:universal-open-side-twisting-cochain} and
  \ref{defn:universal-koszul-obstruction-complex}.  Assume the
  filtrations on \(\Omega C_{\mathrm{op}}\) and \(B_{\mathrm{cl}}\) are
  bounded below, complete, separated, and strongly convergent, with
  finite-dimensional associated graded pieces in every total degree.
  Assume moreover that the local conilpotence condition in
  Definition~\ref{defn:universal-open-side-twisting-cochain} holds and
  continuous dualization is exact on the finite-window towers.
  Then the following conditions are equivalent.
  \begin{enumerate}[label=(A\arabic*)]
    \item The datum is admissible: the generator substitution
      \(c\mapsto\theta,\ u\mapsto O\) is the completed CE/PV identity,
      \(\tau\) is a twisting cochain, and
      \[
        \kappa_\tau\colon
        \Omega\mathrm{Bar}(A_{\mathrm{op}})
        \longrightarrow B_{\mathrm{cl}}
      \]
      is an admissible filtered quasi-isomorphism with the linear symbol
      of Definition~\ref{defn:universal-open-side-twisting-cochain}.
    \item The obstruction equations are exact:
      \[
        \omega_{\mathrm{win}}=0,\qquad
        \mathcal O_{\mathrm{conilp}}=0,\qquad
        \mathcal O_{\mathrm{MC}}(\tau)=0,\qquad
        H^\bullet\!\left(
          \mathcal O_{\mathrm{top}}
          \oplus
          \mathcal O_{\mathrm{cobar}}(\tau)
        \right)=0 .
      \]
  \end{enumerate}
  When these equivalent conditions hold,
  \[
    B_{\mathrm{cl}}
      \simeq
    A_{\mathrm{op}}
  \]
  in the homotopy category of complete augmented filtered
  \(A_\infty\)-algebras.  If the CE coalgebra is also used as the closed
  Koszul source, exact continuous duality identifies the dual form as
  \[
    C^\bullet_{\mathrm{CE},\mathrm{cont}}(\mathfrak g)
      \simeq
    A_{\mathrm{op}}^!
  \]
  precisely when the CE-side cobar map
  \(\Omega C_*^{\mathrm{CE}}(\mathfrak g)\to A_{\mathrm{op}}\) has the
  same exact associated-graded cone and \(\varprojlim^1\) obstruction.
\end{thm}

\begin{proof}
  The equality \(\omega_{\mathrm{win}}=0\) is equivalent to strict
  finite-window compatibility of the CE differential and the Schouten
  differential, with \(P_0\)-pairings added only in the
  kernel-admissible tensor category; this is exactly
  Lemma~\ref{lem:closed-ce-pv-window-limit} and
  Theorem~\ref{thm:universal-formal-local-ce-pv-recognition}.  The
  vanishing of \(\mathcal O_{\mathrm{top}}\) is the Milnor and variance
  condition identifying cohomology of the completed CE/PV objects with
  the admissible finite-window calculation.

  The local conilpotence hypothesis makes the Maurer--Cartan sum finite
  on each finite codomain window, hence a continuous curvature.  The
  equation \(\mathcal O_{\mathrm{MC}}(\tau)=0\) is then the
  Maurer--Cartan equation in the continuous convolution dg Lie algebra.
  It is equivalent to the statement that the adjoint map
  \(\kappa_\tau\) commutes with the completed cobar differential and the
  \(A_\infty\)-operations.  This equality gives the morphism tested
  below.

  After the Maurer--Cartan equation holds, filter
  \(\Omega C_{\mathrm{op}}\) and \(B_{\mathrm{cl}}\) by tensor length
  and by the finite coordinate windows.  The mapping-cone summand in
  \(\mathcal O_{\mathrm{cobar}}(\tau)\) is acyclic exactly when
  \(\operatorname{gr}\kappa_\tau\) is a quasi-isomorphism.  The two
  \(\varprojlim^1\)-summands are exactly the Milnor defects for the
  completed source and target towers.  Strong convergence and
  completeness then make the spectral-sequence comparison theorem
  equivalent to admissible filtered quasi-isomorphism of
  \(\kappa_\tau\).  Necessity follows from the definition of admissible
  filtered quasi-isomorphism in the envelope: it is detected on the
  associated graded and on the exact finite-window towers.

  The bar-cobar counit
  \(\Omega\mathrm{Bar}(A_{\mathrm{op}})\to A_{\mathrm{op}}\) is an
  admissible filtered quasi-isomorphism under the same boundedness and
  convergence hypotheses.  Composing its inverse in the homotopy
  category with \(\kappa_\tau\) gives
  \(A_{\mathrm{op}}\simeq B_{\mathrm{cl}}\).  The final statement is the
  same argument applied to the CE coalgebra
  \(C_*^{\mathrm{CE}}(\mathfrak g)\), followed by exact continuous
  dualization.
\end{proof}

\begin{defn}[Tate cofibrantly generated model structure in the envelope]
  \label{defn:tate-model}
  Define generating cofibrations $I$ and acyclic cofibrations $J$
  by
  $$I=\bigl\{F^w S^{n-1}\hookrightarrow F^w D^n
       \mid n\in\Z_{\geq 0},\ w\in\N\bigr\},\qquad
    J=\bigl\{0\hookrightarrow F^w D^n
       \mid n\in\Z_{\geq 0},\ w\in\N\bigr\},$$
  where $S^{n-1},D^n$ are the standard sphere and disk chain
  complexes in $\Cat{Vec}$ and the Tate window $F^w(-)$ takes the
  $w$-th piece of the filtration. A morphism $f\colon V\to W$ in
  $\Cat{TateChain}_{\geq 0}^{\mathrm{adm}}$ is:
  \begin{enumerate}
    \item a \emph{weak equivalence} if it is an admissible filtered
      quasi-isomorphism in the sense of \alpharef:
      a quasi-isomorphism on every quotient $V/F^wV\to W/F^wW$ and
      on the associated graded;
    \item a \emph{cofibration} if it is in the class generated by
      $I$ under transfinite composition, retracts, and pushouts;
    \item a \emph{fibration} if it has the right lifting property
      against $J$.
  \end{enumerate}
\end{defn}

\begin{thm}[Admissible-envelope model axiom]
  \label{thm:tate-model-structure}
  Under Statement~\ref{stmt:tate-model-envelope}, the category
  \(\Cat{TateChain}_{\geq 0}^{\mathrm{adm}}\) is a cofibrantly generated
  symmetric monoidal model category whose weak equivalences are
  admissible filtered quasi-isomorphisms.  The sets \(I\) and \(J\) of
  Definition~\ref{defn:tate-model} are the chosen generating
  cofibrations and acyclic cofibrations.  The obstruction to replacing
  this axiom by a construction is the simultaneous verification of
  local presentability, smallness, Hovey recognition, the pushout-product
  axiom, the monoid axiom, and compatibility with finite-window
  truncations.
\end{thm}

\begin{proof}
  The assertion is the model-category part of
  Statement~\ref{stmt:tate-model-envelope}.  The strict
  Mittag-Leffler category proves exactness of inverse limits with
  surjective transition maps, but it does not prove smallness for
  transfinite cell complexes, stability of acyclic cofibrations under
  pushout, or the monoid axiom.  Hovey recognition gives the displayed
  model category once those hypotheses are included in the envelope.
\end{proof}

\begin{lem}[Tail quotient obstruction invisible to finite-window weak equivalence]
\label{lem:tate-envelope-tail-boundary}
  Let \(\{V_d\}_{d\geq1}\) be nonzero finite-dimensional vector
  spaces with zero differential, and put
  \[
    V_{\oplus}=\bigoplus_{d\geq1}V_d,\qquad
    V_{\Pi}=\prod_{d\geq1}V_d .
  \]
  Filter both complexes by degree tails.  The inclusion
  \(V_{\oplus}\hookrightarrow V_{\Pi}\) is an admissible filtered weak
  equivalence in the sense of
  Statement~\ref{stmt:tate-model-envelope}: it is an isomorphism on
  every finite quotient and on associated graded pieces.  Its cofiber has
  \[
    H^0(V_{\Pi}/V_{\oplus})=V_{\Pi}/V_{\oplus}\neq0 .
  \]
\end{lem}

\begin{proof}
  The quotient by the tail \(d>w_0\) is
  \[
    V_{\oplus}/\bigoplus_{d>w_0}V_d
      \cong \bigoplus_{d\leq w_0}V_d
      \cong \prod_{d\leq w_0}V_d
      \cong V_{\Pi}/\prod_{d>w_0}V_d .
  \]
  Thus the inclusion is an isomorphism on every finite quotient.  The
  associated graded in degree \(d\) is \(V_d\) on both sides, so the
  associated graded map is also an isomorphism.  This is precisely an
  admissible filtered weak equivalence.

  Since the differential is zero, ordinary cohomology is the underlying
  vector space.  A sequence with one nonzero component in every degree
  defines a nonzero class of \(V_{\Pi}/V_{\oplus}\).  Hence the cofiber
  has nonzero \(H^0\).  An ordinary quasi-isomorphism of chain complexes
  would have acyclic cofiber, so the inclusion is only a filtered weak
  equivalence in the admissible envelope.
\end{proof}


\subsubsection{Operad-algebra transfer}

Let $\mathrm{Lie}_{\Cat{TateVec}}$ be the Lie operad in
$\Cat{TateVec}$, regarded as a $\Sigma$-cofibrant $\Sigma$-module
(by the standard Koszul resolution of the $\Sigma$-cofibrant
replacement), and similarly $\mathrm{Ass}^+_{\Cat{TateVec}}$ for
augmented associative. A $\mathrm{Lie}_{\Cat{TateVec}}$-algebra is
a Tate Lie algebra in $\Cat{TateVec}$; an
$\mathrm{Ass}^+_{\Cat{TateVec}}$-algebra is a Tate complete
augmented associative algebra. Write
$$\mathrm{TateLie}=\mathrm{Lie}_{\Cat{TateVec}}\text{-alg}_{\geq 0},
  \qquad \mathrm{TateAssocAlg}^+=
  \mathrm{Ass}^+_{\Cat{TateVec}}\text{-alg}_{\geq 0}$$
in $\Cat{TateChain}_{\geq 0}^{\mathrm{adm}}$.

\begin{lem}[Transferred model structures on Tate operad-algebras]
  \label{lem:transferred-model}
  Suppose the admissible envelope also satisfies operadic admissibility
  for \(\mathcal P\in\{\mathrm{Lie}_{\Cat{TateVec}},
  \mathrm{Ass}^+_{\Cat{TateVec}}\}\): path objects exist for free
  \(\mathcal P\)-algebras, pushouts along generating acyclic
  cofibrations are weak equivalences, and filtered assembly preserves
  these weak equivalences.  Then the free-forget adjunction
  $\mathrm{Free}_{\mathcal P}\dashv\mathrm{Forget}$ for
  \(\mathcal P\) transfers the cofibrantly
  generated model structure of \ref{thm:tate-model-structure} to
  $\mathcal P\text{-alg}_{\geq 0}$ in the sense of
  Hinich~\cite{hinich-htpy-alg}*{\S 2.2}: the weak equivalences and
  fibrations in $\mathcal P\text{-alg}_{\geq 0}$ are the maps
  whose underlying chain map in
  $\Cat{TateChain}_{\geq 0}^{\mathrm{adm}}$ is one,
  and the cofibrations are determined by lifting.
\end{lem}

\begin{proof}
  Hinich's transfer theorem requires exactly the hypotheses listed above:
  \begin{enumerate}[label=(\arabic*)]
    \item a path object on every algebra of the form
      $\mathrm{Free}_{\mathcal P}(F^w D^n)$;
    \item that pushouts along generating acyclic cofibrations be
      weak equivalences.
  \end{enumerate}
  Item~(1) is the path-object part of operadic admissibility.  Item~(2)
  is the pushout-stability part.  Windowwise finite-dimensional
  arguments are necessary tests, but they do not replace the transfer
  hypothesis in the Tate envelope.
\end{proof}

The dual coalgebra side requires more care. Hinich's transfer
applies to algebra categories; the coalgebra case is governed instead
by conilpotence and the bar-cobar adjunction.
The substitute is the Aubry-Lemaire-Smirnov \cite{loday-vallette}
\S 11.4 model structure on conilpotent dg coalgebras transferred
through the bar-cobar adjunction itself.

\begin{defn}[Conilpotent Tate coalgebras]\label{defn:tate-conilp-coalg}
  Let
  $$\begin{aligned}
  \mathrm{TateConilpCoAlg}
  =\bigl\{\,(C,\Delta_C,d_C)\,:\,&
  C\in\Cat{TateChain}_{\geq 0}^{\mathrm{adm}},\\
  &\Delta_C\colon C\to C\widehat\otimes C
    \text{ continuous coproduct},\\
  &\text{coaugmentation filtration exhaustive}\,\bigr\}.
  \end{aligned}$$
  with morphisms continuous coalgebra maps. Conilpotency means
  the coaugmentation filtration $F^{(n)}_{\mathrm{coaug}}C=
  \ker(\overline\Delta^{(n)})$ is exhaustive: $C=\bigcup_n
  F^{(n)}_{\mathrm{coaug}}C$.
\end{defn}

\begin{lem}[Transferred model structure on Tate conilpotent
  coalgebras]\label{lem:tate-conilp-model}
  Suppose the conilpotent coalgebra transfer of
  Lef\`evre--Hasegawa is valid in the chosen admissible Tate envelope:
  \(\Omega\) preserves coalgebra cofibrations, detects the declared
  weak equivalences, and the unit
  \(C\to\mathrm{Bar}\,\Omega C\) is an admissible filtered weak
  equivalence for every conilpotent Tate coalgebra.  Then
  $\mathrm{TateConilpCoAlg}$ admits a model structure with
  \begin{itemize}
    \item weak equivalences = maps $f\colon C\to C'$ such that
      $\Omega(f)\colon\Omega(C)\to\Omega(C')$ is a weak
      equivalence in $\mathrm{TateAssocAlg}^+$
      (\ref{lem:transferred-model});
    \item cofibrations = monomorphisms of underlying filtered
      chain complexes;
    \item fibrations = right lifting against acyclic cofibrations.
  \end{itemize}
  Equivalently, this is the Lefevre-Hasegawa
  \cite{lefevre-hasegawa}*{Th.~1.3.1.2} transferred model structure
  in the admissible Tate envelope.
\end{lem}

\begin{proof}
  Lefevre-Hasegawa, with proof sketched in Loday-Vallette
  \S 11.4.10, constructs the conilpotent-coalgebra transferred
  model structure in the symmetric monoidal category of $\Q$- or
  $\C$-vector spaces. The proof passes through the cobar-bar
  adjunction $\Omega\colon\mathrm{ConilpCoAlg}\rightleftarrows
  \mathrm{AugAssocAlg}\colon\mathrm{Bar}$ and uses that
  \begin{enumerate}[label=(\roman*)]
    \item $\Omega$ takes coalgebra cofibrations (= mono with
      filtered cokernel) to algebra cofibrations;
    \item $\Omega$ inverts coalgebra weak equivalences;
    \item the unit $\eta_C\colon C\to\mathrm{Bar}\,\Omega(C)$ is a
      weak equivalence for every conilpotent dg coalgebra.
  \end{enumerate}
  Items~(i)--(iii) are precisely the transfer hypotheses in the Tate
  envelope.  Each finite window is a necessary check against the
  ordinary Lef\`evre--Hasegawa proof, but the passage to the completed
  object uses the admissible filtered weak equivalence of the envelope.
\end{proof}

\subsubsection{Bar-cobar comparison in the admissible envelope}

\begin{thm}[Tate bar-cobar comparison in the admissible envelope]
  \label{thm:tate-bar-cobar-quillen}
  Assume Statement~\ref{stmt:tate-model-envelope},
  Lemma~\ref{lem:transferred-model}, and
  Lemma~\ref{lem:tate-conilp-model}, and assume that the bar-cobar unit
  and counit are admissible filtered weak equivalences on the
  conilpotent and augmented objects under consideration.  Then
  The bar-cobar adjunction
  $$\Omega\colon\mathrm{TateConilpCoAlg}\rightleftarrows
    \mathrm{TateAssocAlg}^+\colon\mathrm{Bar}$$
  gives a derived equivalence on the full subcategory generated by
  those admissible cofibrant-fibrant objects.  A Quillen equivalence
  of the entire Tate model categories requires the unit and counit
  weak-equivalence hypotheses for all cofibrant and fibrant objects.
  In particular:
  \begin{enumerate}
    \item $\Omega$ is left Quillen: it preserves cofibrations and
      acyclic cofibrations.
    \item $\mathrm{Bar}$ is right Quillen: it preserves
      fibrations and acyclic fibrations.
    \item For every admissible cofibrant conilpotent coalgebra \(C\)
      and every admissible fibrant augmented associative algebra \(A\)
      in the named subcategories, the unit
      $\eta_C\colon C\to\mathrm{Bar}\,\Omega(C)$ and the counit
      $\varepsilon_A\colon\Omega\,\mathrm{Bar}(A)\to A$ are weak
      equivalences.
    \item The induced derived adjunction restricts to an equivalence
      on the homotopy subcategories generated by the admissible
      objects satisfying the displayed unit and counit hypotheses.
  \end{enumerate}
\end{thm}

\begin{proof}
  Items~(1) and~(2) are immediate from
  \ref{lem:tate-conilp-model}: weak equivalences in
  $\mathrm{TateConilpCoAlg}$ are defined as maps whose $\Omega$
  is a weak equivalence, and fibrations on the coalgebra side are
  defined by lifting.

  For item~(3), the unit $\eta_C$: filter both sides by symmetric
  / tensor length on the underlying $\Cat{TateVec}$ object, plus
  the Tate window filtration $F^w$. The associated graded gives,
  windowwise, a finite-dimensional dg coalgebra and the standard
  bar-cobar resolution
  $C^{\mathrm{gr}}\to\mathrm{Bar}\,\Omega(C^{\mathrm{gr}})$. By
  Loday-Vallette \cite{loday-vallette}*{\S 2.2.5, \S 11.1.5}, this
  is a quasi-isomorphism on each window. Conilpotency
  (\ref{defn:tate-conilp-coalg}) ensures the coaugmentation
  filtration on $C$ is exhaustive and Hausdorff, so the abutment
  of the spectral sequence of the filtration computes the global
  qiso. Mittag-Leffler exactness in the Tate filtration
  (\alpharef\ item~1) lifts the windowwise qiso to a global
  filtered qiso.

  For the counit
  \(\varepsilon_A\colon\Omega\,\mathrm{Bar}(A)\to A\), filter both
  sides by augmentation/tensor length and by the Tate window.  On every
  finite window and every finite length cap this is the ordinary
  bar-cobar counit for an augmented dg associative algebra over
  \(\C\), hence a quasi-isomorphism by the standard normalized bar
  resolution \cite{loday-vallette}*{\S 2.2.5, \S 11.1.5}.  The
  filtrations are bounded below, complete, and Hausdorff in the
  admissible envelope, so the spectral sequence comparison and
  Mittag-Leffler exactness lift the windowwise quasi-isomorphism to the
  completed Tate object.

  In the Lie specialization \(A=\widehat U(\mathfrak g)\), the same
  counit can also be read through the PBW filtration:
  \(\mathrm{gr}_{\mathrm{PBW}}\widehat U(\mathfrak g)\cong
  \widehat S_{\mathrm{Tate}}(\mathfrak g)\), and the associated graded
  comparison is the symmetric Koszul resolution of
  Loday-Vallette \cite{loday-vallette}*{\S 3.4.13}.  This is the
  closed-side PBW form used later, but the Quillen counit above is the
  associative bar-cobar counit for general \(A\).

  Item~(4) is formal from~(1)-(3) on the stated admissible
  cofibrant-fibrant subcategory.  On the whole model categories,
  Quillen's adjoint-functor criterion is the assertion that the same
  unit and counit are weak equivalences for every cofibrant and fibrant
  object.
\end{proof}

\subsubsection[Promotion to the associated infinity-categories]{Promotion to the associated $\infty$-categories}

\begin{cor}[Infinity-categorical Tate bar-cobar comparison]
  \label{cor:tate-bar-cobar-infinity}
  Apply the Dwyer-Kan / Hammock-localization functor to the
  admissible subcategory comparison of \ref{thm:tate-bar-cobar-quillen}:
  $$N_\Delta\bigl(\mathrm{TateConilpCoAlg}^{\mathrm{adm,cf}},W\bigr)
    \simeq
    N_\Delta\bigl(\mathrm{TateAssocAlg}^{+,\mathrm{adm,cf}},W\bigr)$$
  as the $\infty$-categories associated to the admissible model
  subcategories in the sense of Lurie's
  \emph{Higher Algebra} \S A.3, where $\mathrm{cf}$ denotes the
  cofibrant-fibrant objects satisfying the unit and counit
  hypotheses, and \(W\) is the inherited class of weak equivalences.
  The equivalence is induced by \(L\Omega\) and \(R\,\mathrm{Bar}\) on
  this admissible subcategory.
\end{cor}

\begin{proof}
  Theorem~\ref{thm:tate-bar-cobar-quillen} gives inverse derived
  functors on the admissible cofibrant-fibrant subcategories.  The
  Dwyer--Kan localization of a relative equivalence of such
  subcategories is an equivalence of the associated
  \(\infty\)-categories.
\end{proof}

\subsubsection{Identification with the Lie operadic Koszul duality}

The bar-cobar comparison of \ref{thm:tate-bar-cobar-quillen} is
\emph{associative} bar-cobar.  Specialization to the Lie operad
gives the cochain-level comparison used in the CE/PV theorem.

\begin{cor}[Tate Lie bar-cobar / CE Koszul duality]
  \label{cor:tate-lie-bar-cobar}
  For $\mathfrak g$ a lower-central Tate-pronilpotent graded Lie algebra in the
  sense of \alpharef, the admissible bar-cobar comparison applied
  to $A=\widehat U(\mathfrak g)\in\mathrm{TateAssocAlg}^+$ yields
  the equivalence in the homotopy category
  $$L\Omega\bigl(C_*^{\mathrm{CE}}(\mathfrak g)\bigr)
    \xrightarrow{\sim}\widehat U(\mathfrak g),\qquad
    R\,\mathrm{Bar}\bigl(\widehat U(\mathfrak g)\bigr)
    \xrightarrow{\sim}C_*^{\mathrm{CE}}(\mathfrak g),$$
  and hence, after continuous dualization, a Lie-Koszul-dual equivalence
  of cochain algebras
  $$C^\bullet_{\mathrm{CE},\mathrm{cont}}(\mathfrak g)
    \simeq\bigl(\widehat U(\mathfrak g)\bigr)^!$$
  in the corresponding homotopy category of complete augmented
  Tate algebras.

	  No symmetric monoidal $\infty$-categorical lift is asserted from
	  this corollary.  Such a lift requires a weak monoidal Quillen
	  equivalence for the chosen Tate operadic model and a separate
	  proof that the Lie-operadic base change preserves the completed
	  Koszul duality.
\end{cor}

\begin{proof}
  Apply \ref{thm:tate-bar-cobar-quillen} to the
  Lie-operadic Koszul-dual pair $(\mathrm{Lie},\mathrm{CoComm})$
  in $\Cat{TateVec}$. Loday-Vallette \cite{loday-vallette}*{Th.~13.4.6}
  identifies the operadic bar-cobar with the Lie-bar-cobar of
  CE chains and completed enveloping algebra.  The statement is
  therefore an equivalence in the homotopy category determined by the
  admissible envelope.  Monoidal functoriality is a further structure,
  not a consequence of the associative bar-cobar counit.
\end{proof}

\begin{lem}[Closed-side Tate formalism]
\label{lem:closed-tate-koszul-formalism}
  Let \(\mathfrak h=\varprojlim_M\mathfrak h_{\leq M}\) be a closed
  Hamiltonian Tate Lie algebra whose finite windows are
  cochain-CE/PV-admissible, and kernel-admissible when a \(P_0\) kernel is
  used, in the sense of
  Definition~\ref{defn:closed-window-obstruction-terms}.  Assume:
  \begin{enumerate}[label=(K\arabic*)]
    \item the shifted cotangent extension
      \(\mathfrak g=\mathfrak h\ltimes D_{\mathrm{cont}}[1]\) is an
      object of \(\Cat{TateChain}_{\geq0}^{\mathrm{adm}}\);
    \item the CE coalgebra \(C_*^{\mathrm{CE}}(\mathfrak g)\) is
      conilpotent for the coaugmentation filtration;
    \item the completed enveloping algebra
      \(\widehat U(\mathfrak g)=\varprojlim_M\widehat U(\mathfrak g_{\leq M})\)
      is complete and separated for the PBW filtration;
    \item the finite-window PBW associated graded maps assemble to
      \[
        \mathrm{gr}_{\mathrm{PBW}}\widehat U(\mathfrak g)
          \cong
        \widehat{\mathbf S}_{\mathrm{Tate}}(\mathfrak g).
      \]
  \end{enumerate}
  Then the closed-side formalism has two compatible consequences:
  \[
    C^\bullet_{\mathrm{CE},\mathrm{cont}}(\mathfrak g)
      \xrightarrow{\ \Phi\ }
    \mathrm{PV}\bigl(\widehat{\mathbf S}_{\mathrm{Tate}}(\mathfrak h)\bigr)
  \]
  is the admissible completed CE/PV identity obtained from the
  finite-window family of Lemma~\ref{lem:closed-ce-pv-window-limit}
  together with the exact continuous-dualization and variance data in
  (K1)--(K4), and
  \[
    L\Omega C_*^{\mathrm{CE}}(\mathfrak g)
      \simeq \widehat U(\mathfrak g),\qquad
    R\,\mathrm{Bar}\,\widehat U(\mathfrak g)
      \simeq C_*^{\mathrm{CE}}(\mathfrak g)
  \]
  is the Tate Lie bar-cobar / CE Koszul duality of
  Corollary~\ref{cor:tate-lie-bar-cobar}.  The two constructions are
  compatible with the same finite-window transition maps.
\end{lem}

\begin{proof}
  The first assertion is Lemma~\ref{lem:closed-ce-pv-window-limit} plus the
  exact continuous-dualization and variance hypotheses included above.
  The obstruction terms \(\mathfrak o^{\mathrm{br}}\) and
  \(\mathfrak o^{\mathrm{coad}}\) are the cochain CE/PV transition
  obstructions; \(\mathfrak o^{\mathrm{Cas}}\) is the additional
  kernel/\(P_0\)-pairing obstruction.  Their vanishing gives the
  finite-window compatibility; the admissible envelope gives the
  cochain inverse-limit datum for
  \(\Phi^{\mathrm{CE/PV}}_{\mathrm{cont}}
    =\varprojlim_M\Phi_M^{\mathrm{CE/PV}}\).

  For the Koszul assertion, hypotheses (K1)--(K4) put
  \(C_*^{\mathrm{CE}}(\mathfrak g)\) and
  \(\widehat U(\mathfrak g)\) in the admissible model envelope with
  complete filtrations.  Conilpotency is the coalgebra hypothesis
  needed for \(\mathrm{Bar}\,\Omega\), and PBW completeness identifies
  the associated graded of the enveloping algebra with the completed
  symmetric algebra.  Corollary~\ref{cor:tate-lie-bar-cobar} therefore
  applies.  Both constructions are obtained by first applying the
  finite-dimensional CE/PV and bar-cobar statements on
  \(\mathfrak g_{\leq M}\), then passing through the same
  Mittag-Leffler inverse limit; hence the transition maps are the same.
\end{proof}

\subsubsection{Module reconstruction from the CE/PV identity}

\begin{thm}[Admissible-envelope module reconstruction from the CE/PV identity]
\label{thm:admissible-envelope-module-reconstruction}
  Let $\mathfrak h_{\mathrm{adm}}$ be a pro-finite-dimensional Lie algebra in
  $\Cat{TateVec}$ whose shifted-cotangent extension
  $\mathfrak g_{\mathrm{adm}}
  =\mathfrak h_{\mathrm{adm}}\ltimes
  \mathfrak h_{\mathrm{adm},\mathrm{cont}}^\vee[1]$
  satisfies the lower-central Tate-pronilpotent admissible-envelope hypotheses of
  \ref{thm:tate-bar-cobar-quillen}.  Assume further that:
  \begin{enumerate}[label=(M\arabic*)]
    \item the strict cochain map
  $$\Phi\colon C^\bullet_{\mathrm{CE},\mathrm{cont}}(\mathfrak g_{\mathrm{adm}})
    \longrightarrow \mathrm{PV}\bigl(
    \widehat S_{\mathrm{Tate}}(\mathfrak h_{\mathrm{adm}})\bigr)$$
      is the admissible CE/PV quasi-isomorphism obtained from the
      finite-window family of Lemma~\ref{lem:closed-ce-pv-window-limit}
      after the required variance and exact continuous-dualization data
      are given;
    \item the completed dg \(P_0\)-algebras on both sides lie in an
      admissible transferred \(P_0\)-algebra model structure inside
      \(\Cat{TateChain}_{\geq0}^{\mathrm{adm}}\), with weak
      equivalences detected on underlying admissible filtered
      quasi-isomorphisms;
    \item the module categories are formed in the same locally
      presentable monoidal envelope, and the two algebras are replaced by
      cofibrant-fibrant representatives when needed.
  \end{enumerate}
  Then \(\Phi\) induces an equivalence, inside that admissible Tate
  model envelope, of module $\infty$-categories over the two identified
  completed dg \(P_0\)-algebras
  $$\widetilde\Phi\colon
    N_\Delta\bigl(C^\bullet_{\mathrm{CE},\mathrm{cont}}(\mathfrak g_{\mathrm{adm}})
      \text{-mod},W\bigr)
    \xrightarrow{\sim}
    N_\Delta\bigl(\mathrm{PV}\bigl(\widehat S_{\mathrm{Tate}}(\mathfrak h_{\mathrm{adm}})
      \bigr)\text{-mod},W\bigr)$$
  in the \(\infty\)-categorical setting of Lurie's \emph{Higher Algebra}
  \S 5.5.  The
  hypotheses (M1)--(M3) presuppose an already constructed admissible
  \(P_0\)-algebra envelope. The conclusion is the algebraic
  $\infty$-categorical equivalence inside that envelope. The completed
  dg \(P_0\) model structure and the factorization-centre theorem for
  the full untruncated formal Hamiltonian algebra are separate
  constructions.
\end{thm}

\begin{proof}
  Assumption (M1) gives a strict admissible filtered quasi-isomorphism of
  completed dg \(P_0\)-algebras.  The strictness is the generator-level
  identity \(c\mapsto\theta,\ u\mapsto O\), together with
  \(d_\pi\Phi=\Phi d_{\mathrm{CE}}\) and bracket compatibility from the
  finite-window CE/PV calculation.

  By (M2), this strict weak equivalence is a weak equivalence in the
  transferred \(P_0\)-algebra model structure.  Replace the two
  \(P_0\)-algebras by cofibrant-fibrant representatives if necessary.
  Extension and restriction of scalars along a weak equivalence of
  admissible algebra objects give a Quillen equivalence of module model
  categories; this is the standard Schwede--Shipley/Lurie mechanism
  for module categories over weakly equivalent algebra objects
  \cite{schwede-shipley-monoid} and
  \cite{lurie-ha}*{Th.~7.1.4.6}.  Applying
  Dwyer--Kan localization gives the displayed equivalence of module
  \(\infty\)-categories.

  Functoriality in the Lie algebra gives interval-wise compatibility for
  admissible interval objects
  \(\mathfrak h_{\mathrm{adm},I}=
  \Omega^\bullet_c(I)\otimes\mathfrak h_{\mathrm{adm}}\) when these
  interval objects remain in the same admissible envelope.  This gives
	  interval-wise categorical compatibility. Weiss descent and unreduced
	  open-BV centrality are governed by their own descent and centrality
	  criteria.
\end{proof}

\begin{rmk}[Relation to Lurie's higher-algebraic Koszul
  duality]
  \label{rmk:lurie-koszul}
  Lurie \cite{lurie-ha}*{\S 5.5.5} constructs a higher-algebraic
  Koszul duality between augmented $\infty$-operad-algebras and
  their Koszul-dual coalgebras.  For the operad
  $\mathrm{Lie}_{\Cat{TateVec}}$, Koszul self-duality (with
  shift) gives the equivalence of \ref{cor:tate-lie-bar-cobar}
  in the lower-central Tate-pronilpotent setting.  The two
  constructions agree: weak equivalences in the model-categorical
  presentation are admissible filtered quasi-isomorphisms
  (\alpharef), and Lurie A.3.7 with the Hinich transfer of
  \ref{lem:transferred-model} identifies them with the
  Hammock-localized weak equivalences of the $\infty$-categorical
  presentation.
\end{rmk}

\subsubsection{Conilpotence and Casimir convergence}

\begin{prop}[Bar-cobar conilpotence separated from Casimir convergence]
  \label{prop:quillen-vs-casimir}
  The admissible bar-cobar comparison of \ref{thm:tate-bar-cobar-quillen}
  and the resulting $\infty$-categorical equivalence of
  \ref{thm:admissible-envelope-module-reconstruction}, when their
  admissible-envelope hypotheses hold, are independent
  of the convergence of the Tate Casimir element. The cochain
  complexes $C^\bullet_{\mathrm{CE},\mathrm{cont}}(\mathfrak g)$
  and $\mathrm{PV}\bigl(\widehat S_{\mathrm{Tate}}(\mathfrak h)
  \bigr)$ are equivalent in the model category of
  \ref{lem:transferred-model} after imposing the same conilpotent
  admissibility hypotheses. The required hypothesis here is
  conilpotence; Casimir convergence belongs to the analytic kernel
  problem.
\end{prop}

\begin{proof}
  The bar-cobar comparison uses
  $\widehat S^c(\mathfrak g[1])$ and $T(\mathfrak g)$ as
  underlying graded objects. Both are well-defined and conilpotent
  in the admissible envelope under \alpharef\ items~1, 2 (continuous
  dual exact, conilpotent CE coalgebra). The formal Casimir
  \(\sum_d e_{d,i}\otimes e_d^{i,\vee}\) is a kernel hypothesis only when
  it converges in the chosen coefficient category; in the strict
  product/direct-sum Tate endpoint it is precisely the obstruction
  isolated in Lemma~\ref{lem:tate-casimir-obstruction}.  Such a kernel is
  required for the Costello \cite{costello-renormalization} heat-kernel
  propagator construction at the analytic factorization-algebra level
  where the bracket is not strict.  The algebraic
	  bar-cobar comparison uses only the conilpotent admissibility
	  hypotheses: the cochain-level $\Phi$ extends, in the admissible model
	  envelope, to an $\infty$-categorical equivalence.
\end{proof}

\begin{prop}[Strict product/direct-sum kernel obstruction]
\label{prop:strict-product-direct-sum-kernel-obstruction}
  Let \(\mathcal C_{\mathrm{ker}}\) be the category whose objects are
  complete coefficient pairs \((H,D)\), equipped with finite quotients
  \((H_{\leq M},D_{\leq M})\), compatible transition maps, and a tensor
  \(K\in H\widehat\otimes D\) whose image in
  \(H_{\leq M}\otimes D_{\leq M}\) is the finite Casimir \(C_M\) for
  every \(M\).  Morphisms preserve the finite quotients and carry \(K\)
  to \(K'\).  The original strict product/direct-sum pair
  \[
    H_\Pi=\prod_{d\geq1}H_d,\qquad
    D_\oplus=\bigoplus_{d\geq1}H_d^\vee
  \]
  lies outside \(\mathcal C_{\mathrm{ker}}\).  More strongly, any object
  of \(\mathcal C_{\mathrm{ker}}\) mapping to the strict pair by a
  finite-window identity morphism produces the obstructed
  infinite-support tensor in \(H_\Pi\widehat\otimes D_\oplus\).

  Thus initial, terminal, and representing universal properties for
  kernel-admissible completions pass through a genuine kernel object.  A
  weighted product completion is a constructed regulator for
  tail-sensitive ordinary cohomology classes.
\end{prop}

\begin{proof}
  If the strict pair were an object, it would contain a tensor
  \(K\in H_\Pi\widehat\otimes D_\oplus\) whose image in every finite
  window is \(C_M\).  If a kernel-admissible object mapped to the strict
  pair by a finite-window identity morphism, the image of its kernel
  would give the same obstructed tensor.  But \(D_\oplus\) has finite
  support in the degree direction, while the compatible family
  \((C_M)_M\) has a nonzero \(H_d^\vee\)-component for every \(d\).
  This is exactly the support obstruction of
  Theorem~\ref{thm:strict-unweighted-obstruction}.  Therefore the strict
  pair lies outside \(\mathcal C_{\mathrm{ker}}\), and universal-property
  assertions must pass through a kernel-admissible completion.
\end{proof}

\begin{prop}[Residual Weiss-cosheaf condition]
  \label{prop:weiss-cosheaf-residual}
  The $\infty$-categorical equivalence of
  \ref{thm:admissible-envelope-module-reconstruction} on stalks and on
  the locally-constant interval-factorization $\mathcal F_{\mathrm{cl}}$
  of \ref{thm:hamiltonian-current-factorization} gives the cochain-level
  categorical comparison.  The Weiss-cosheaf condition for the
  \emph{unreduced} factorization algebra
  $\mathrm{Obs}^{\mathrm{cl}}_{\mathrm{closed},H}$ on $\R^2\times\C^2$
  requires the additional descent check below.
  The Weiss / descent condition for an unreduced factorization
	  algebra on a higher-dimensional manifold requires checking
	  \v{C}ech descent for Weiss covers, an independent
	  cosheaf-theoretic statement.  Costello--Gwilliam Vol.~I
	  \S\S 6.1, 6.5, 7.2 provide the factorization, cosheaf, and
	  basis-extension vocabulary; the present argument uses them as
	  vocabulary and records the unreduced descent obligation as a residual
	  class.  The
  admissible Tate
  Quillen equivalence closes the cochain-level
  $\infty$-categorical comparison; the Weiss descent for the unreduced
  factorization algebra remains a separate cosheaf-theoretic
  obligation requiring the heat-kernel propagator extension of
  Problem~\ref{prob:formal-factorization-center}.  The two-dimensional
  formal singular product calculation is the separate result of
  Theorem~\ref{thm:formal-two-dimensional-holomorphic-ope}.
\end{prop}

\begin{proof}
  Lurie \cite{lurie-ha}*{\S 5.5.4}, with the same terminology as
  Costello-Gwilliam \cite{costello-gwilliam}*{Vol.~I \S 6.4}: a locally constant factorization
  algebra on $\R$ determines an $E_1$-algebra up to the standard
  contractible choice; the
  bar-cobar comparison of \ref{thm:tate-bar-cobar-quillen}
  realizes this on the admissible subcategory.  On a higher-dimensional manifold, the Weiss
  cover $\{U_\alpha\}$ has structure-map descent independent of
  the algebraic bar-cobar; in particular the heat-kernel
  propagator on $\R^2\times\C^2$ enters the analytic check.
  This heat-kernel propagator is the analytic datum relating the
  bar-cobar identity to Weiss descent.
\end{proof}
